% HBL §2.13 supplementary mathematical calculation; NO new Lean acceptance.
\subsection*{A Mellin far-transport inverse and the non-gaining splice commutator}
For the original frozen Hartree--Kummer radial operator put
\[
\mathcal L_{\mu,\theta}=\partial_r^2+
 \left(\frac2r+\frac{ir}{2}\right)\partial_r+
 \theta+i\left(\frac14-\mu\right),\qquad
 A_{\mu,\theta}=\frac i2r\partial_r+\beta,\quad
 \beta=\theta+i\left(\frac14-\mu\right).
\]
In the physical radial Hilbert space
\(\mathscr H=L^2((0,\infty),4\pi r^2\,dr)\), the unitary logarithmic
transform \(\mathsf Mf(x)=\sqrt{4\pi}e^{3x/2}f(e^x)\) gives
\[
 \mathsf MA_{\mu,\theta}\mathsf M^{-1}
 =\frac i2\partial_x+\theta-i(\mu+\tfrac12).
\]
Fourier transformation in \(x\) therefore gives
\[
 \|A_{\mu,\theta}^{-1}F\|_{\mathscr H}
 \le\frac{\|F\|_{\mathscr H}}{\mu+1/2},
 \qquad \mu>-\frac12.
\]
The resolvent identity yields, uniformly for \(\mu_j+1/2\ge\eta_*>0\),
\[
\begin{aligned}
\|A_{\mu_1,\theta_1}^{-1}F_1-A_{\mu_2,\theta_2}^{-1}F_2\|_{\mathscr H}
&\le\eta_*^{-1}\|F_1-F_2\|_{\mathscr H}\\
&\quad+\eta_*^{-2}(|\mu_1-\mu_2|+|\theta_1-\theta_2|)
\|F_2\|_{\mathscr H}.
\end{aligned}
\]
For the \emph{original untruncated} inverse-square mixed term
\[
 \mathcal C(z)=|x|^{-2}*
 \bigl(|x|^{-1}\operatorname{Re}z\,|z|^2\bigr),
 \qquad z\in\dot H^1_{\rm rad}(\mathbb R^3),
\]
Hardy, Sobolev and HLS give
\[
 \|A_{\mu,\theta}^{-1}\mathcal C(z)\|_{\mathscr H}
 \le\frac{2C_{\rm HLS}S_3^2}{\mu+1/2}\|\nabla z\|_2^3.
\]
This gives a genuine infinite-dimensional \emph{far principal-part}
response, unlike a two-speed finite-rank correction. It is not a
full Kummer inverse and has no asserted renormalized-time decay.

There are two exact obstructions to inserting this far inverse
globally. For a smooth real radial radiation snapshot with \(z(0)\ne0\),
angular integration of the original kernel gives
\[
 \mathcal C(z)(r)=4\pi z(0)^3\log(1/r)+O(1)\qquad(r\downarrow0).
\]
If \(\mu>1/4\) and \(p=2(\mu-1/4)+2i\theta=2i\beta\),
the unique radial \(L^2\) inverse can be written
\[
 G(r)=2i\int_r^\infty(r/\rho)^p\mathcal C(z)(\rho)
 \,\frac{d\rho}{\rho}
 =\frac{4\pi z(0)^3}{\beta}\log(1/r)+O(1).
\]
Hence \(G\notin H^2(B_1)\); the actual radial Laplacian cannot
be applied to it in global \(L^2\) merely from its far inverse bound.

Moreover, for \(y=\ell r\), \((\mathsf U_\ell g)(r)=
\ell^{3/2}g(\ell r)\), and a smooth cutoff \(\chi(y)\),
\[
 \mathcal L_{\mu,\theta}\mathsf U_\ell g
 =\mathsf U_\ell(A_{\mu,\theta}+\ell^2\Delta_{\rm rad})g,
\]
\[
 \boxed{
 [A_{\mu,\theta}+\ell^2\Delta_{\rm rad},\chi]g
 =\frac{iy}{2}\chi'g+
 \ell^2\left(2\chi'g'+(\chi''+2\chi'/y)g\right).
 }
\]
The first interface commutator carries \emph{no} small factor
\(\ell\). Thus a generic cutoff does not produce a positive-time
residual gain; this order-\(b\) term must be handled by the
\emph{same actual} selected Fredholm/Cramer/Green interface and the
complete Hartree PDE. The smooth-radiation calculation is a
no-shortcut obstruction, not a counterexample to an unknown selected
HBL trajectory.

The positive \(\mu\) bound, the moving source and parameter
derivatives, the full nonlocal four-region residual, its positive
time order and Lipschitz control, and the finite physical recursion
are all still open. This calculation is \emph{paper only};
the revised original formal acceptance remains \(11/13\).
